\documentclass[11pt]{article}
\usepackage{amsmath,amssymb,amsthm,hyperref}
\newtheorem{theorem}{Theorem}
\newtheorem{lemma}[theorem]{Lemma}
\newcommand{\Q}{\mathbb Q}
\newcommand{\Ann}{\operatorname{Ann}}
\title{Nonlinear relative GST folds: complete pure-model classification and secondary memory}
\author{Hand derivation before Lean; author self-review}
\date{3 November 2026}
\begin{document}\maketitle
\section{Owner, inputs and output category}
Checkpoint: \texttt{2e0324f1328409d68c25d85897435187d0a380b0}.
Keep arbitrary original smooth projective complex $X$, $p\geq1$ and
prescribed rational Hodge $\alpha$ as the geometric input. The universal
cycle target is still a rational cycle on this same $X$ with canonical
Betti class $\alpha$. The results here concern the actual rational torsor
model and its marked base, not a replacement definition of Chow classes.

For connected $X$ let $A=\bigoplus_q\operatorname{Hdg}^q(X)_\Q$, with
cohomological degrees $2q$, augmentation ideal $\mathfrak m=A_{>0}$ and
$\mathfrak m^{\dim X+1}=0$. For any finite homogeneous curvature list
$a_i\in A_{p_i}$ set
\[
 B(a)=A\otimes\Lambda(u_i,v_i),\quad |u_i|=|v_i|=2p_i-1,
 \quad du_i=a_i,\quad dv_i=0,\qquad I(a)=(a_i).
\]
All maps called relative fix $A$ pointwise. Arbitrary nonlinear graded
maps are allowed; their images need not be linear in the handles. A
zigzag below consists of quasi-isomorphisms of DGAs under $A$.
These are unfiltered equivalences unless a filtration is explicitly named.

\section{The full base kernel, without a linear-fold restriction}
\begin{theorem}[Exact marked-base kernel]
The kernel of the canonical map $A\longrightarrow H^*(B(a))$ is exactly
$I(a)$. Consequently $A/I(a)$ is the actual image subalgebra of the base
in cohomology. It need not be the whole cohomology algebra.
\end{theorem}
\begin{proof}
Expand an arbitrary element uniquely by exterior word length. The
length-zero component of its differential is $\sum_i b_i a_i$, where
$b_i$ is its coefficient of the one-letter word $u_i$. Longer words
cannot contribute length zero. Therefore a base element which is a
boundary lies in $I(a)$. Conversely $\sum_i b_i a_i=d(\sum_i b_i u_i)$,
since all coefficients in $A$ are even and closed. This proves equality,
including an empty curvature list.
\end{proof}

For the full model $R\otimes\Lambda(u_i,v_i)$ with
$R=H^*(X,\Q)$, including odd and off-diagonal base classes, the identical
argument gives kernel $I(a)R$, with the usual graded sign in the coefficient
of $u_i$. Its contraction to $A$ is exactly $I(a)$. Here is why this last
assertion does not assume a rational projection onto complex Hodge types.
In each degree the multiplication map
$\bigoplus_i H^{2(q-p_i)}(X,\Q)(-p_i)\to H^{2q}(X,\Q)$ is a morphism of
polarizable rational pure Hodge structures. The rational polarizing
orthogonal complement of its kernel is a Hodge substructure: the kernel
is stable under the Weil operator, so positivity makes the restricted
polarization nondegenerate. This complement maps isomorphically to the
image. A rational Hodge target in that image therefore has a rational
Hodge preimage, supplying coefficients in $A_{q-p_i}$. This proves the
contraction statement and keeps all original base cohomology.

\begin{theorem}[Nonlinear and quasi-isomorphic invariance]
For any DGA $C$ under $A$, the kernel of $A\to H(C)$ is preserved by every
relative quasi-isomorphism, hence every relative zigzag. If $C$ has a
chain retraction to $(A,0)$, that kernel is zero. Thus nonzero $I(a)$
cannot be erased by any relative nonlinear equivalence, even into a
presentation whose differentials are not pure base curvatures.
\end{theorem}
\begin{proof}
The cohomology map of a relative quasi-isomorphism is an isomorphism
commuting with the marked base map; their kernels coincide. A chain
retraction makes the base map split injective on cohomology. In particular,
if every target handle differential has positive handle length, setting
all handles to zero is such a retraction. No inspection of the nonlinear
coordinate formula is needed.
\end{proof}

\section{Complete classification of the pure paired presentations}
\begin{theorem}[Classification in the larger category]
For two finite pure paired models over this fixed $A$, the following are
equivalent:
\begin{enumerate}
\item they are joined by a relative DGA quasi-isomorphism zigzag;
\item $I(a)=I(b)$ as graded ideals and their multiplicities
$r_q=\#\{i:p_i=q\}$ agree in every weight;
\item they are related by a homogeneous $A$-linear change of the paired
handle coordinates, hence by a relative DGA isomorphism.
\end{enumerate}
Therefore the previous minimum $\mu_q=\dim_\Q(I/\mathfrak m I)_q$ remains
sharp among all pure paired presentations in the entire nonlinear relative
quasi-isomorphism class. The equivalence need not preserve an ordered
basis. This does not classify all DGAs with nonlinear differentials.
\end{theorem}
\begin{proof}
The base-kernel theorem recovers $I$ from (1). To recover the handle
multiplicities, take the derived fiber along $A\to\Q$. The model is a
finite semifree $A$-module, filtered by exterior length, with differential
lowering length. It is also a relative free DGA obtained by attaching the
odd generators with differential already in $A$. Its ordinary fiber
therefore computes the derived fiber:
\[
 \Q\otimes_A^{\mathbb L}B(a)\simeq
 (\Lambda_\Q(u_i,v_i),0).
\]
For clarity the flatness assertion is not ordinary tensor preservation
for arbitrary modules. A finite semifree complex is homotopy projective:
lifting a map into an acyclic complex and its nullhomotopy is performed
successively over its free filtration. Thus a quasi-isomorphism between
such models has acyclic homotopy-projective cone and a chain homotopy
inverse. Tensoring that homotopy with $\Q$ preserves it. For an arbitrary
zigzag use semifree replacements at each intermediate object; the same
lifting argument defines the invariant derived fiber. These constructions
also preserve the product in its derived algebra. The graded
indecomposables of its cohomology exterior algebra have dimension $2r_q$
in degree $2q-1$. This recovers every $r_q$.

For (2) implies (3), apply the preceding graded-local normal-form theorem
to both lists. Their minimal generators have the same counts $\mu_q$.
Writing one minimal list in terms of the other gives a homogeneous matrix
whose equal-weight diagonal blocks are invertible over $\Q$, since both
lists are bases of $I/\mathfrak m I$. Its off-diagonal part is strictly
weight triangular and has a finite inverse. Their closed pairs have equal
counts $r_q-\mu_q$ and can be matched. Apply the same matrix to each
conjugate partner list. This gives the claimed isomorphism; it also
preserves the specified model filtration. (3) implies (1) immediately.
\end{proof}

The marking matters. On $A=\Q[x,y]/(x^2,y^2)$, $|x|=|y|=2$, swapping
$x,y$ identifies the absolute models of curvatures $x$ and $y$. Their
marked kernels $(x)$ and $(y)$ differ, so they are not equivalent under
the original base. Total Betti numbers alone cannot make this distinction.

\section{An actual nonlinear fold and its exact transported differential}
On $X=\mathbb P^1\times\mathbb P^1$, use curvatures $(x,0,xy)$ with
handle pairs $(u,v),(z,t),(w,s)$ of degrees $(1,1),(1,1),(3,3)$.
Thus $du=x$, $dw=xy$ and all other generator differentials vanish. For
every $\lambda\in\Q$ the degree-preserving substitution
\[
 U=u,\quad V=v,\quad Z=z,\quad T=t,\quad
 W=w-yu+\lambda uvz,\quad S=s
\]
is an actual nonlinear graded algebra isomorphism. Its inverse is
$w=W+yu-\lambda uvz$. Transporting the full differential gives
\[
 dW=\lambda xvz,
\]
not zero when $\lambda\ne0$. The transformed model still has base
kernel $(x)$. The explicit closed coordinate is $W-\lambda uvz$.
Consequently the same complete closed-pair factor is recovered, without
mistaking the absence of the base term of $dW$ for closedness. Every
monomial, both inverse orders, all absolute degrees and the nonlinear
Leibniz signs are checked in the accompanying finite calculation.

\section{Complete one-curvature cohomology ring and secondary memory}
\begin{theorem}[One paired handle, including multiplication]
For any even graded commutative $A$ and homogeneous $a$ of degree $2p$,
\[
 H\bigl(A\otimes\Lambda(u,v),du=a,dv=0\bigr)
 \cong \bigl(A/(a)\ltimes (\Ann_A(a)u)\bigr)\otimes\Lambda(v).
\]
The summand containing $u$ is a square-zero ideal, shifted by $2p-1$.
The action of $A/(a)$ on it is well-defined because $a\Ann_A(a)=0$.
\end{theorem}
\begin{proof}
In the two-term $u$ factor, $b+cu$ is closed exactly when $ca=0$;
all boundaries are the base elements $aA$. Multiplying two $u$ terms
gives zero. Changing a base representative by $ab$ changes its product
with $cu$ by $abc u=0$. This proves the algebra formula, not just a
vector-space count. The flat $v$ factor is an exterior algebra with zero
differential, so tensoring preserves the formula with its graded signs.
\end{proof}

This ring formula does not imply formality. Take the actual rational
torsor over $\mathbb P^1\times\mathbb P^1$ with $a=xy$, $p=2$. The
classes $[x][y]$ and $[y]^2$ vanish. With primitives $u$ and zero,
the triple Massey product $\langle[x],[y],[y]\rangle$ contains
$[-uy]$ (the opposite global sign convention is equivalent here).
It is closed because $d(uy)=xy^2=0$, and is nonzero because there are no
$u$-component boundaries. The complete indeterminacy is
$[x]H^3+H^3[y]=\operatorname{span}_\Q\{[xv],[yv]\}$, since
$H^3=\Q[v]$. The class $[yu]$ is independent of that span, so zero is
not in the product. This gives an explicit higher operation retained by
the actual torsor; the ring is not being used as a false zero-differential
replacement. This fixture tests the arbitrary-base structural formulas;
it is not the universal Chow target.

\section{Geometric disposition and proof boundary}
For the original arbitrary prescribed $(X,p,\alpha)$, construct its same
rational homotopy-fiber torsor and any finite cup births. Its full base
kernel is $(\alpha)H^*(X,\Q)$, with rational Hodge contraction $(\alpha)$.
That is the exact computed original-base effect under every relative
nonlinear quasi-isomorphism, and the pure paired presentations now have
a complete classification. The one-handle ring also computes its full
primary multiplicative memory and identifies possible secondary memory.
No algebraicity of $\alpha$ was needed for these topological constructions.
They do not construct a universal Chow cycle or prove the Hodge conjecture.

The derivation and classification are hand proofs with author self-review.
Lean additions encode the base-boundary submodule, exact first-coefficient
calculation, relative chain-equivalence invariance, retraction obstruction,
and concrete two-curvature packets. They do not claim a formalized derived
category, Massey-product construction or arbitrary-X geometric comparison.
Compilation is left to the owner. No historical-priority claim is made.
\end{document}
